We conduct 960 additional runs covering 48 distinct settings, each using the same 20 seed labels as the primary study. Table~\ref{tab:sensitivity-design} gives the reference values and the alternative test values. Seven primary settings provide the references: no organization, full organization, social-only organization, the near and far sites at 30\% shelter capacity, and one and five boats at the far site with 10\% shelter capacity. Warning lead time, hospital capacity and people per boat trip are varied across all seven. Depth reduction is varied only at full organization; the social-capital increase is varied only at full and social-only organization. We omit parameter changes that do not apply to a setting. Settings unaffected by a parameter change reuse the verified primary runs with the same seeds and do not count as new runs.

We examine four changes: from no to full organization, from no to social-only organization, from the near to the far shelter at 30\% capacity, and from one to five boats at the far site with 10\% capacity. For each tested parameter value, we calculate the mortality difference within each seed and average across seeds. We then subtract that seed's primary mortality difference before averaging again. This second quantity, a paired difference of differences, measures how much the parameter change alters the comparison. Reusing the same seeds preserves the dependence created by shared reference runs.

Exact sign-flip tests use the integer death-count changes and the same sign-exchangeability assumption as the primary tests. Holm correction covers all 30 applicable changes in comparisons. Each approximate 95\% Monte Carlo interval based on Student's t distribution concerns one estimate, not all estimates jointly. The design was fixed after primary analysis and before the sensitivity runs. It tests selected parameter values individually; it does not measure each parameter's share of total outcome variation or the consequences of changing several parameters together.
